%!TEX root =  tfg.tex
\chapter{Contexto}

\begin{quotation}[Novelist]{Ernest Hemingway (1899--1961)}
The good parts of a book may be only something a writer is lucky enough to overhear or it may be the wreck of his whole damn life -- and one is as good as the other.
\end{quotation}

\begin{abstract}
En este apartado se describe el marco general del proyecto, 
analizando la relevancia de la industria del videojuego, 
la evolución del género roguelike y cómo la obra clásica de Dante Alighieri 
sirve como base conceptual para la propuesta técnica y artística de Via Inferni.
\end{abstract}

\section{El mundo de los videojuegos y el género Roguelike}

La industria del videojuego se ha consolidado como uno de los sectores de entretenimiento más dinámicos a nivel tecnológico. Dentro de este mercado, el género roguelike destaca por su enfoque en la rejugabilidad y el desafío. Se caracteriza por tres pilares que se aplican en este proyecto: la generación aleatoria de escenarios (para que ninguna partida sea igual a la anterior), la muerte permanente (donde el progreso se reinicia al morir) y el combate táctico en tiempo real.

\section{El Infierno de Dante como base narrativa}

El proyecto se enmarca en la adaptación de la literatura clásica al entorno interactivo. El juego utilizar la estructura de los nueve círculos del infierno descritos por Dante Alighieri:
\begin{itemize}
\item \textbf{Estructura:} Cada círculo actúa como un nivel con ambientación y enemigos específicos, desde el Limbo hasta la traición en el noveno círculo.
\item \textbf{Personajes:} La dualidad entre Dante (fuerza) y Virgilio (habilidades) se traslada a mecánicas de juego diferenciadas para aportar profundidad estratégica.
\end{itemize}

\section{Desarrollo Técnico y Metodología}

Desde el subcontexto de la Ingeniería del Software, el proyecto se desarrolla utilizando herramientas profesionales para garantizar la calidad del producto:
\begin{itemize}
\item \textbf{Motor:} Unity 2D para la lógica de juego y renderizado.
\item \textbf{Gestión:} Uso de Git, GitHub Project y GitHub para el control de versiones, asegurando un flujo de trabajo colaborativo organizado.
\end{itemize}

\section{Estado del Arte}

Actualmente, los videojuegos independientes (indies) que utilizan técnicas similares gozan de gran popularidad. Tecnológicamente, el uso de algoritmos que generen conexiones de salas en rejillas (grid) representa una solución eficiente para crear contenido extenso con recursos optimizados. Via Inferni se posiciona en este mercado como una propuesta que une la rigurosidad técnica de la ingeniería con la riqueza cultural de una obra universal.
